\begin{textsample}{POS Dim 1 – sp – Score 28.00 – sp/sp\_ef\_ciencias\_humanas\_geo\_2.txt}  \label{ex:f1_pos_164}
O foco do ensino de Geografia \textbf{hoje} está no estudo do espaço geográfico, conceito \textbf{que} pode ser entendido como produto das relações sociais, econômicas, políticas, culturais, simbólicas e ambientais \textbf{que} nele se estabelecem.
Nessa perspectiva, as relações definidas entre os elementos naturais e os construídos pela atividade humana, são regulados pelo “ tempo da natureza ” ( processos bioquímicos e físicos, responsáveis pela produção e interação dos objetos naturais ) e pelo “ tempo histórico ” ( \textbf{marcas} acumuladas pela atividade humana como produtora de artefatos sociais ).
O espaço geográfico ainda pode ser entendido como resultado da trama entre objetos técnicos e informacionais, fluxos de matéria e informação, \textbf{que} se manifestam e atuam sobre \textbf{uma} base física.
Para Santos ( 2008 ), a natureza do espaço é a soma do resultado material acumulado das ações humanas através do tempo e, de outro, animado pelas ações \textbf{atuais} \textbf{que} lhe atribuem \textbf{um} dinamismo e \textbf{uma} funcionalidade.
A paisagem tem sido tomada como \textbf{um} primeiro foco de análise, como ponto de partida para aproximação de seu objeto de estudo \textbf{que} é o espaço geográfico.
Pode ser definida como a unidade visível do real e \textbf{que} incorpora todos os fatores resultantes da construção natural, social e cultural.
Para Santos ( 1997 ), a paisagem pressupõe, também, \textbf{um} conjunto de formas e funções em constante transformação, seus aspectos “ visíveis ”, mas, por outro \textbf{lado}, as formas e as funções indicam a estrutura espacial, em princípio, “ invisível ”, e resulta sempre do casamento da paisagem com a sociedade.
Já para Vitte ( 2007 ), o conceito de paisagem se manifesta como polissêmico e resultado de \textbf{uma} representação filosófica e social ; cada sociedade, por meio de sua cultura, imprime \textbf{uma} particular plasticidade à natureza \textbf{que} é produzida pela intencionalidade social.
Já para Ab’Saber ( 2003 ), as paisagens têm sempre o caráter de herança de processos ( fisiográficos e biológicos ), de atuação antiga, remodelados e modificados por processos de atuação recente.
São \textbf{uma} herança, \textbf{um} patrimônio coletivo dos povos \textbf{que}, \textbf{historicamente}, os modificaram ao longo do tempo e do espaço.
A definição de lugar está cada vez mais complexa, global e dinâmica.
O lugar pode ser entendido como o espaço \textbf{que} se torna próximo do indivíduo, constituindo -se como o lugar do pertencimento, encontros, experiência, dimensão afetiva, identidade, subjetividade e lugar do simbólico.
No contexto \textbf{atual}, a sociedade depara -se com \textbf{um} conjunto de acontecimentos \textbf{que} ultrapassam as \textbf{fronteiras} do local, pois são eventos globais, mas sua repercussão se materializa no lugar.
Aliás, o lugar é o depositário final dos eventos, de acordo com Santos ( 2003 ).
Ainda para o \textbf{autor} ( 2008 ), o lugar abarca \textbf{uma} permanente mudança, decorrente da própria lógica da sociedade e das inovações técnicas \textbf{que} estão sempre transformando o espaço geográfico.
Com relação ao território, pode ser considerado sinônimo de espaço \textbf{vivido}, apropriado, usado, delimitado, \textbf{que} \textbf{configura} os aspectos políticos, econômicos, ambientais e culturais.
O território não é apenas a configuração política de \textbf{um} Estado-Nação, mas sim o espaço construído pela formação social.
Segundo Raffestin ( 1993 ), o território não poderia ser nada mais \textbf{que} o produto dos \textbf{atores} sociais.
São eles \textbf{que} produzem o território, partindo da realidade inicial dada, \textbf{que} é o espaço.
Ainda para o \textbf{autor}, o território é definido com base em \textbf{um} sistema composto por nós e redes, \textbf{que} constrói \textbf{uma} estrutura \textbf{conceitual}, como limite, \textbf{fronteiras}, vizinhança, \textbf{territorialidade}, entre outros.
Já para Haesbaert ( 2007 ), o território é sempre múltiplo, diverso, complexo e imerso em relações de dominação e/ou de apropriação sociedade-espaço, desdobra -se da dominação político-econômica mais concreta e \textbf{funcional} à apropriação mais subjetiva e/ou cultural-simbólica.
Segundo Corrêa ( 1998 ), o conceito de região, tradicionalmente, é entendido como \textbf{uma} parte da superfície da Terra, dimensionada segundo escalas territoriais diversificadas, caraterizada pelos elementos da natureza ou como \textbf{uma} paisagem e sua extensão territorial, na qual se entrelaçam os componentes humanos e a natureza.
Ao longo da história, o conceito foi reformulado e está associado à ideia de território amplo, regionalização, \textbf{divisão} do espaço, localização, extensão de \textbf{um} fenômeno, entre outros.
Outro conceito estruturante refere -se à educação \textbf{cartográfica}, \textbf{visto} \textbf{que} a linguagem \textbf{cartográfica} tem \textbf{um} papel importante no processo de aprendizagem em Geografia, no sentido de contribuir para o desenvolvimento de habilidades necessárias para o entendimento das interações, dinâmicas, relações e dos fenômenos geográficos em diferentes escalas e para a formação da cidadania e da criticidade e autonomia do estudante.
A cartografia escolar vem se estabelecendo como \textbf{um} conhecimento construído nas \textbf{interfaces} entre Cartografia, Educação e Geografia.
No entanto, a cartografia escolar abrange conhecimentos e práticas para o ensino de conteúdos originados na própria cartografia, mas \textbf{que} se caracteriza por lançar mão de visões de diversas áreas.
Em seu estado \textbf{atual}, pode referir -se a formas de se apresentar conteúdos relativos ao espaço-tempo social, a concepções \textbf{teóricas} de diferentes áreas de conhecimento a ela relacionadas, a experiências em diversos contextos culturais e a práticas com tecnologias da informação e comunicação. ( ALMEIDA, 2011, p.07 ) Para Castellar ( 2005 ), a cartografia é considerada \textbf{uma} linguagem, \textbf{um} sistema de código de comunicação \textbf{imprescindível} em todas as esferas da aprendizagem em Geografia, articulando fatos e conceitos.
Ressalta -se \textbf{que} também pode ser entendida como técnica e pode se tornar \textbf{uma} metodologia inovadora, na medida em \textbf{que} permite relacionar conteúdos, conceitos e fatos.
As pesquisas desenvolvidas pela autora ( 2011 e 2017 ) \textbf{revelam} \textbf{que} a alfabetização \textbf{cartográfica}, ao ensinar a ler em Geografia, cria condições para \textbf{que} o estudante leia.
Esse processo de alfabetização \textbf{cartográfica} ocorre de forma gradual, em função da complexidade das relações, dinâmicas e dos fenômenos estudados, da faixa etária do estudante e da necessidade de construção de referenciais espaciais.
Na infância, o estudante experimenta o grafismo como forma de expressão e o desenho pode ser considerado \textbf{uma} das primeiras manifestações do processo de alfabetização.
Em seguida, com \textbf{um} repertório ampliado, representa cartograficamente o espaço, tendo como base elementos presentes no seu lugar de vivência.
Desse modo, ao reconhecer os elementos \textbf{constituintes} do espaço e as \textbf{inter-relações} com outros espaços, o estudante amplia o seu repertório \textbf{conceitual} e \textbf{metodológico}, construindo os conhecimentos geográficos e \textbf{cartográficos} no decorrer do Ensino Fundamental e, posteriormente, no Ensino Médio.
As tecnologias no ensino de Geografia apresentam formas de observar o espaço em diversas escalas, subsidiando a compreensão das relações ambientais, sociais, econômicas, políticas e culturais em diferentes tempos.
As Geotecnologias \textbf{revelam} potencial \textbf{didático-pedagógico} e têm possibilitado cada vez mais \textbf{que} o estudante tenha acesso a diferentes dados e representações gráficas e \textbf{cartográficas} produzidas pelo Sensoriamento Remoto, por Sistemas de Informações Geográficas ( SIG ), pelo Sistema de \textbf{Posicionamento} Global ( GPS ) e pela Cartografia Digital.
Nesse conjunto de possibilidades para o fortalecimento do ensino de Geografia no Ensino Fundamental, destaca -se a contribuição da Cartografia Inclusiva para o processo de aprendizagem dos estudantes.
Carmo e Sena ( 2018 ) em suas pesquisas apontam \textbf{que} os princípios da cartografia tátil \textbf{que}, originalmente, foram pensados para estudantes com deficiência visual, mas \textbf{que}, com o uso nas salas regulares, se mostraram \textbf{interessantes} para todos os estudantes.
Considerando os pontos destacados, a educação \textbf{cartográfica} contribui para a educação para a cidadania, por meio de \textbf{uma} aprendizagem significativa, contextualizada e inclusiva, em \textbf{que} os estudantes mobilizam diversas competências, habilidades e conhecimentos para ler e interpretar o espaço geográfico.

% matched lemmas: ator, atual, autor, cartográfico, conceitual, configurar, constituinte, didático-pedagógico, divisão, fronteira, funcional, historicamente, hoje, imprescindível, inter-relação, interessante, interface, lado, marca, metodológico, posicionamento, que, revelar, territorialidade, teórico, um, ver, viver
\end{textsample}
